\chapter{Introduction}
\label{cha:intro}

In today’s interconnected world, traditional rule-based Intrusion Detection Systems (IDS) struggle to keep pace with the rapid evolution of cyber threats. Attackers continuously adapt their methods by exploiting trust relationships, user behaviour, and system misconfigurations, rendering static detection techniques increasingly ineffective. Machine Learning (ML) has therefore been widely proposed as a suitable approach to address these challenges, and numerous benchmark studies demonstrate its effectiveness in this domain. For example, Abrar et al. show, using the well-known NSL-KDD dataset, that classifiers such as Random Forest (RF), Extra Tree, Decision Tree, and Support Vector Machines (SVM) can achieve very high detection accuracy for classifying normal versus malicious traffic, often exceeding 99\% \cite{abrar2020machine}. Furthermore, Gawand and Meesala confirm that NSL-KDD remains a useful benchmark for prototyping and comparative evaluation, despite the increasing complexity of real-world network traffic \cite{gawand2025improved}. At the same time, NSL-KDD is a benchmark dataset and does not fully capture the characteristics of production networks. Hence, results obtained on this dataset should therefore be interpreted as comparative and methodological rather than indicative of real-world performance.

The need for adaptive intrusion detection is further reinforced by recent large-scale cyber incidents. In 2025, Ingram Micro suffered a ransomware attack by the SafePay group that resulted in the exfiltration of over 3.5 TB of data through compromised VPN credentials. Similarly, Orange SA was targeted by the Warlock group, while Louis Vuitton (LVMH) experienced a data breach via a third-party CRM vulnerability, exposing customer information across multiple European regions. These incidents share common characteristics, such as third-party security vulnerabilities, credential theft, and social-engineering-based intrusions, which frequently bypass traditional IDS solutions. This highlights the practical relevance of ML-based intrusion detection for both research and real-world deployment\cite{pkware2024breaches}.

In this project, intrusion detection is formulated as a supervised classification task using the NSL-KDD dataset, in which preprocessed network connection records are provided as input to machine-learning models that predict whether network traffic is normal or malicious and, in extended experiments, identify the corresponding attack category (i.e. multi-class classification), including Denial of Service (DoS), Probe, Remote-to-Local (R2L), and User-to-Root (U2R) attacks. This work builds upon prior research that has extensively evaluated ML techniques on NSL-KDD. Tavallaee et al. introduced NSL-KDD as a cleaned and rebalanced alternative to the original KDD Cup ’99 dataset, addressing issues such as redundant records and biased attack distributions that previously led to inflated performance estimates \cite{tavallaee2009detailed}. Subsequent studies consistently report that tree-based ensemble models outperform many alternatives due to their ability to capture nonlinear feature interactions and handle class imbalance effectively \cite{abrar2020machine, hegde2024intrusion, hong2021machine}. Other authors have shown that, while deep learning models can achieve strong performance, classical machine-learning methods such as Random Forest often match or exceed their results while offering lower computational cost and greater interpretability \cite{ali2025deep}.

Guided by these findings, this project evaluates multiple machine-learning models for intrusion detection on NSL-KDD, starting with Logistic Regression as a linear baseline, followed by a Random Forest classifier and finally a neural network to assess whether increased model expressiveness could yield additional benefits. By systematically comparing these approaches, the project aims to provide both empirical insight and practical guidance on the suitability of different ML techniques for network intrusion detection.